\documentclass[../main.tex]{subfiles}

\begin{document}

\chapter{Inequalities}

Inequality is perhaps the most amiable, but extremely in depth
section to study in Algebra. It could range from simple inequalities
such as solving for $x$ in $x + 1 > 2$ to problems involving
numerous advanced techniques and known results. As a result,
although few inequalities in advanced problems could be proven by
bashing, most would require an elegant solution. With that in mind,
this chapter will begin with numerous classical inequalities such as
the AM-GM inequality to advanced concepts such as majorization.

\section{Classic Inequalities}

One of the first concepts that most students, including myself,
study when getting started with olympiad is I believe inequalities.
Many problems are solved by tweaking known inequalities and doing
some clever substitutions. In this note, we will discuss fundamental
and classic inequalities as an introduction to olympiad inequalities.

\subsection{AM-GM Inequality}

Let's start with AM-GM Inequality. I assume most of you know or
at least have heard of this inequality as we learn it in middle and
high school. AM-GM inequality stands for Arithmetic Mean-Geometric
Mean inequality and literally compares those two means.

\begin{definition}
{AM-GM Inequality}
{AM-GM Inequality}
For $a_k \geq 0$, the following inequality holds with equality if
and only if $a_1 = a_2 = \dots = a_n$.
\[
\frac{ \sum_{i=1}^n a_i }{n} \geq
\sqrt[n]{ \prod_{i=1}^n a_i }
\]
\end{definition}

You might be more familiar with $a + b \geq 2 \sqrt{ab}$, but the
inequality above is the generalized version.

\begin{tip}
AM-GM as a fundamental inequality has lots of application. However
in olympiads, it is used the most for cancellation and investigating
the minimum.
\end{tip}

Consider the following example.

\begin{problem}{}{}
Prove the following inequality for $a, b, c, d \geq 0$.
\[
(a + b)(b + c)(c + d)(d + a) \geq 16abcd
\]
\end{problem}
\begin{solution}
Because $a, b, c, d$ are non-negative, we could utilize AM-GM
inequality. By AM-GM inequality, the following inequalities hold.
\begin{align*}
    a + b &\geq 2 \sqrt{ab} \\
    b + c &\geq 2 \sqrt{bc} \\
    c + d &\geq 2 \sqrt{cd} \\
    d + a &\geq 2 \sqrt{da}
\end{align*}
Multiplying the inequalities,
\[
(a + b)(b + c)(c + d)(d + a) \geq 
2 \sqrt{ab} \cdot 2 \sqrt{bc} \cdot 
2 \sqrt{cd} \cdot 2 \sqrt{da} = 16abcd
\]
hold and $(a + b)(b + c)(c + d)(d + a) \geq 16abcd$.
\end{solution}

\begin{problem}{}{}
SECOND PROBLEM FOR AM-GM INEQUALITY
\end{problem}

When we discuss arithmetic and geometric mean, we normally
include weighted means. Thus, it is natural to consider AM-GM
inequality with weighted means.

\subsubsection{Weighted AM-GM Inequality}

Let's first get started with definitions.

\begin{definition}
{Weighted AM-GM Inequality}
{Weighted AM-GM Inequality}
For weights $w_1, w_2, \dots, w_n \geq 0$ such that
$\sum_{i=1}^n w_i = w$ the following inequality holds.
\[
\frac{\sum_{i=1}^n w_i a_i}{w}
\geq \sqrt[w]{\prod_{i=1}^n a_i^{w_i}}
\]
The equality holds if and only if $a_1 = a_2 = \dots = a_n$.
\end{definition}

From the definition, we can notice that the original inequality is
obtained with $w_1 = w_2 = \cdots = w_n = \frac{1}{n}$. Here is
a problem on weighted AM-GM inequality.

\begin{problem}
{2020 IMO Problem 2}
{2020 IMO Problem 2}
Show that for $a, b, c, d \in \mathbb{R}$ such that
$a \geq b \geq c \geq d > 0$ and $a + b + c + d = 1$,
\[
(a + 2b + 3c + 4d) a^a b^b c^c d^d < 1
\]
holds.
\end{problem}

(\href{https://youtu.be/Vz9BC5q7zdc}{Video Solution})

\begin{solution}
First, consider the following weighted AM-GM inequality for
$a, b, c, d$ with weights $a, b, c, d$ respectively.
\begin{align*}
    \frac{a \cdot a + b \cdot b + c \cdot c + d \cdot d}
        {a + b + c + d} 
    &\geq \sqrt[a + b + c + d]{a^a b^b c^c d^d} \\
    a^2 + b^2 + c^2 + d^2
    &\geq a^a b^b c^c d^d
\end{align*}
Therefore, it suffices to show that the following inequality holds.
\[
(a + 2b + 3c + 4d)(a^2 + b^2 + c^2 + d^2) < 1
\]
Consider inequalities below.
\begin{align*}
    a^2 (a + 2b + 3c + 4d)
    &< a^2 (a + 3b + 3c + 3d) \quad (\because b \geq d) \\
    b^2 (a + 2b + 3c + 4d)
    &< b^2 (3a + b + 3c + 3d) \quad (\because 2a \geq b + d) \\
    c^2 (a + 2b + 3c + 4d)
    &< c^2 (3a + 3b + c + 3d) \quad (\because 2a + b \geq 2c + d) \\
    d^2 (a + 2b + 3c + 4d)
    &< d^2 (3a + 3b + 3c + d) \quad (\because 2a + b \geq 3d)
\end{align*}
Integrating,
\begin{align*}
    (a + 2b + 3c + 4d)(a^2 + b^2 + c^2 + d^2)
    &< a^2 (a + 3b + 3c + 3d) + b^2 (3a + b + 3c + 3d) \\
    &\qquad + c^2 (3a + 3b + c + 3d) + d^2 (3a + 3b + 3c + d) \\
    &\qquad < (a + b + c + d)^3 = 1
\end{align*}
and the inequality holds.
\end{solution}

With the ideas of weighted AM-GM inequality in mind, we can further
generalize our results for different means.

\subsection{Power Mean and Weighted Power Mean Inequality}

Power Mean inequality, also known as QM-AM-GM-HM inequality,
compares the four major means with generalization of $AM-GM$
inequality. Without simple notations, we can represent the
Power Mean inequality as $M_2 \geq M_1 \geq M_0 \geq M_{-1}$.

\begin{definition}
{Power Mean Inequality}
{Power Mean Inequality}
For positive real numbers $a_i$, the inequality
\[
\sqrt{ \frac{\sum_{i=1}^n a_i^2}{n} }
\geq \frac{\sum_{i=1}^n a_i}{n}
\geq \sqrt[n]{\prod_{i=1}^n a_i}
\geq \frac{n}{ \sum_{i=1}^n \frac{1}{a_i} }
\]
holds with equality if and only if $a_1 = \cdots = a_n$.
\end{definition}

With the definition in mind, we can naturally apply the same for
weighted power means.

\begin{definition}
{Weighted Power Mean Inequality}
{Weighted Power Mean Inequality}
For positive real numbers $a_i$, the inequality
\[
\sqrt{ \frac{\sum_{i=1}^n w_i a_i^2}{w} }
\geq \frac{\sum_{i=1}^n w_i a_i}{w}
\geq \sqrt[w]{\prod_{i=1}^n a_i^{w_i}}
\geq \frac{w}{ \sum_{i=1}^n \frac{w_i}{a_i} }
\]
holds with equality if and only if $a_1 = \cdots = a_n$.
\end{definition}

Continuing from AM-GM inequality and its variations, the second most
taught inequality in middle and high schools is probably
Cauchy-Schwarz Inequality.

\subsection{Cauchy-Schwarz Inequality}

As always, let's get started with the definition.

\begin{definition}
{Cauchy-Schwarz Inequality}
{Cauchy-Schwarz Inequality}
For all real numbers $a_i$ and $b_i$,
\[
\left( \sum_{i=1}^n a_i^2 \right) \left( \sum_{i=1}^n b_i^2 \right)
\geq
\left( \sum_{i=1}^n a_i b_i \right)^2
\]
with equality if and only if $\frac{a_i}{b_i} = k$ for all integer
$1 \leq i \leq n$, where $k > 0$ and $a_ib_i \neq 0$.
\end{definition}

When $n = 2$ we get the form $(a^2 + b^2)(x^2 + y^2) \geq (ax + by)^2$
that we initially learn. Continuing, we can obtain an interesting
direct consequence. The lemma below known as Titu's lemma, T2 Lemma,
Sedrakyan's inequality, and Engel's form and was named after one of
the greatest mathematician Titu Andreescu.

\begin{corollary}
{Titu's Lemma}
{Titu's Lemma}
For positive real numbers $a_i$ and $b_i$,
\[
\sum_{i=1}^n \frac{a_i^2}{b_i}
\geq
\frac{ \left( \sum_{i=1}^n a_i \right)^2 }{ \sum_{i=1}^n b_i }
\]
\end{corollary}
\begin{proof}
Substituting $x_i = \frac{a_i}{\sqrt{b_i}}$ and $y_i = \sqrt{b_i}$
to the Cauchy-Schwarz inequality, the following inequalities are
obtained.
\begin{align*}
    \left( \sum_{i=1}^n x_i^2 \right)
    \left( \sum_{i=1}^n y_i^2 \right)
    &\geq
    \left( \sum_{i=1}^n x_i y_i \right)^2 \\[0.5em]
    \left(
        \sum_{i=1}^n \left( \frac{a_i}{\sqrt{b_i}} \right)^2
    \right)
    \left(
        \sum_{i=1}^n \left( \sqrt{b_i} \right)^2
    \right)
    &\geq
    \left(
        \sum_{i=1}^n \frac{a_i}{\sqrt{b_i}} \cdot \sqrt{b_i}
    \right)^2
\end{align*}
Recasting,
\[
\left( \sum_{i=1}^n \frac{a_i^2}{b_i} \right)
\left( \sum_{i=1}^n b_i \right)
\geq
\left( \sum_{i=1}^n a_i \right)^2
\]
and
\[
\sum_{i=1}^n \frac{a_i^2}{b_i}
\geq
\frac{ \left( \sum_{i=1}^n a_i \right)^2 }{ \sum_{i=1}^n b_i }
\]
is obtained. Thus, the lemma holds.
\end{proof}

This lemma comes especially handy if an inequality involves
fractions with square numbers. Now that we discussed a direct
consequence of Cauchy-Schwarz inequality, we can consider its
generalization.

\subsection{Hölder's Inequality}

Cauchy-Schwarz inequality is actually a special case of Hölder's
inequality. Both inequalities can be used in multiple ways such as
eliminating radicals and fractions. Below is the definition of
Hölder's inequality.

\begin{definition}
{Hölder's Inequality}
{Hölder's Inequality}
Let $a_i, b_i, \ldots, z_i > 0$ and $\lambda_a, \lambda_b, \ldots,
\lambda_z \geq 0$ such that $\lambda_a + \lambda_b + \dots +
\lambda_z = 1$. Then,
\begin{align*}
\left( \sum_{i=1}^n a_i \right)^{\lambda_a}
\left( \sum_{i=1}^n b_i \right)^{\lambda_b} \cdots
\left( \sum_{i=1}^n z_i \right)^{\lambda_z}
\geq
    \sum_{i=1}^n a_i^{\lambda_a} b_i^{\lambda_b}
        \cdots z_i^{\lambda_z}
\end{align*}
holds with equality if $a_1 : a_2 : \dots : a_n \equiv b_1 : b_2 :
\dots : b_n \equiv \dots \equiv z_1 : z_2 : \dots : z_n$.
\end{definition}

Indeed, Cauchy-Schwarz inequality does resemble Hölder's Inequality!
To derive Cauchy-Scharz inequality from Hölder's, it suffices to show
that a specific case of Hölder's inequality will lead to
Cauchy-Schwarz inequality.

Let $\lambda_a = \lambda_b = \frac{1}{2}$. Substituting, the following
inequalities hold.
\begin{align*}
    \left( \sum_{i=1}^n a_i \right)^\frac{1}{2}
    \left( \sum_{i=1}^n b_i \right)^\frac{1}{2}
    \left( \sum_{i=1}^n c_i \right)^0 \cdots
    \left( \sum_{i=1}^n z_i \right)^0
    &\geq
    \sum_{i=1}^n
        a_i^{\frac{1}{2}} b_i^{\frac{1}{2}} \cdot
        c_i^0 \cdots z_i^0 \\
    \left( \sum_{i=1}^n a_i \right)^\frac{1}{2}
    \left( \sum_{i=1}^n b_i \right)^\frac{1}{2}
    &\geq
    \sum_{i=1}^n a_i^\frac{1}{2} b_i^\frac{1}{2}
\end{align*}
Let $\alpha_i = \sqrt{a_i}$ and $\beta_i = \sqrt{b_i}$. Substituting,
\[
\left( \sum_{i=1}^n \alpha_i^2 \right)^\frac{1}{2}
\left( \sum_{i=1}^n \beta_i^2 \right)^\frac{1}{2}
\geq
\sum_{i=1}^n \alpha_i \beta_i
\]
and
\[
\left( \sum_{i=1}^n \alpha_i^2 \right)
\left( \sum_{i=1}^n \beta_i^2 \right)
\geq
\left( \sum_{i=1}^n \alpha_i \beta_i \right)^2
\]
is obtained. This is Cauchy-Schwarz inequality! Below is a very
classic problem on Hölder's Inequality

\begin{problem}{}{}
Show that the following inequality holds for all $a, b, c > 0$.
\[
(a^3 + 2)(b^3 + 2)(c^3 + 2) \geq (a + b + c)^3
\]
\end{problem}

(\href{https://youtu.be/FUtgo54gq7I}{Video Solution})

\begin{solution}
The key to this problem is changing the left-hand side to apply
Hölder's Inequality. Consider the following inequality by Hölder.
\[
(a^3 + 1 + 1)^\frac{1}{3} (1 + b^3 + 1)^\frac{1}{3}
    (1 + 1 + c^3)^\frac{1}{3}
\geq a + b + c
\]
Thus, cubing both sides leads to the desired result.
\end{solution}

Continuing from Cauchy-Schwarz inequality and its variation, another
fundamental inequality of different type is rearrangement inequality.

\subsection{Rearrangement Inequality}

Let's start with definition.

\begin{definition}
{Rearrangement Inequality}
{Rearrangement Inequality}
Let $a_i$ and $b_i$ be real numbers such that $a_1 \geq a_2 \geq
\cdots \geq a_n$ and $b_1 \geq b_2 \geq \cdots \geq b_n$. Moreover,
let $b_{\sigma(i)}$ be the $i^\text{th}$ term of the permutations of
$b_1, \cdots, b_n$. Then, the following inequalities hold.
\[
\sum_{i=1}^n a_i b_i
\geq \sum_{i=1}^n a_i b_{\sigma(i)}
\geq \sum_{i=1}^n a_i b_{n-i+1}
\]
The equality holds if at least one sequence is constant.
\end{definition}

With the definition, we can obtain a special consequence.

\begin{corollary}
{Chebyshev's Inequality}
{Chebyshev's Inequality}
For real numbers $a_1 \geq a_2 \geq \cdots \geq a_n$ and $b_1 \geq b_2
\geq \cdots \geq b_n$,
\[
n \left( \sum_{i=1}^n a_i b_i \right)
\geq \left( \sum_{i=1}^n a_i \right)
\left( \sum_{i=1}^n b_i \right)
\geq n \left( \sum_{i=1}^n a_i b_{n-i+1} \right)
\]
holds with the equality if at least one sequence is constant.
\end{corollary}
\begin{proof}
Consider the following cases where $b_1 \geq b_2 \geq \cdots \geq b_n$
are cycled through in the rearrangement inequality.
\begin{align*}
    \sum_{i=1}^n a_i b_i
    &\geq a_1 b_1 + a_2 b_2 + \cdots + a_n b_n
    \geq \sum_{i=1}^n a_i b_{n-i+1} \\
    \sum_{i=1}^n a_i b_i
    &\geq a_1 b_2 + a_2 b_3 + \cdots + a_n b_1
    \geq \sum_{i=1}^n a_i b_{n-i+1} \\
    &\qquad\qquad\qquad \vdots \\
    \sum_{i=1}^n a_i b_i
    &\geq a_1 b_n + a_2 b_1 + \cdots + a_n b_{n-1}
    \geq \sum_{i=1}^n a_i b_{n-i+1} \\
\end{align*}
Adding the inequalities,
\[
n \left( \sum_{i=1}^n a_i b_i \right)
\geq \left( \sum_{i=1}^n a_i \right)
\left( \sum_{i=1}^n b_i \right)
\geq n \left( \sum_{i=1}^n a_i b_{n-i+1} \right)
\]
and the corollary holds.
\end{proof}

\end{document}

